Sprint 1 xây dựng nền tảng kỹ thuật, quản lý tài liệu bằng Git và hiện thực MVP xem menu bằng \textbf{Next.js, NestJS và PostgreSQL}. Kết quả nổi bật: \textbf{2 API sản phẩm, 8 sản phẩm mẫu, 18/18 test đạt} trong lần kiểm local ngày \textbf{23/10/2026}; đã có minh chứng CI thành công và phát hiện lỗi.

\section{Thiết lập repository}

Dự án được tổ chức tại \href{https://github.com/DuVinhLoi/BrewLite}{DuVinhLoi/BrewLite}. Mỗi ứng dụng có manifest và lockfile riêng, cài dependency bằng \code{npm ci}. Node.js được thống nhất ở phiên bản \textbf{24.19.0}.

\begin{table}[H]
\caption{Cấu trúc repository và quy trình làm việc}
\small
\begin{tabularx}{\linewidth}{L{0.25\linewidth}Y}
\toprule
\thead\textbf{Thành phần} & \textbf{Tạo phẩm và vai trò} \\
\midrule
Frontend và backend & \code{apps/web}: giao diện Next.js; \code{apps/api}: NestJS, Prisma và PostgreSQL. \\
Tài liệu và cấu hình & \code{docs/}, README, \code{.gitignore}, \code{.env.example} và cấu hình mẫu theo app; không đưa env thật vào Git. \\
Branch và review & Nhánh \code{a2/Task4}, \href{https://github.com/DuVinhLoi/BrewLite/pull/9}{PR \#9} hướng vào \code{develop-Assignment2}; quy trình yêu cầu ít nhất \textbf{1 reviewer} khác tác giả. \\
Quy ước và DoD & Conventional Commits; hoàn thành tiêu chí chấp nhận, kiểm tra chất lượng, cập nhật tài liệu và truy vết, review và xử lý blocker. \\
Product Backlog & Có backlog cho \textbf{2 Sprint}: Sprint 1 về menu và nền tảng; Sprint 2 về đặt hàng, thanh toán và bàn giao. \\
\bottomrule
\end{tabularx}
\end{table}

Nguồn: \gitfile{README.md}{README}, \gitfile{CONTRIBUTING.md}{CONTRIBUTING} và \gitfile{docs/sprint/backlog.md}{Product Backlog}. Hồ sơ chưa có phân công cá nhân, lịch Sprint và URL board trực tuyến; đề cho phép sử dụng Product Backlog lưu trong repository.

\section{Quản lý tài liệu có kiểm soát phiên bản}

Tài liệu được lưu trong \code{docs/} và quản lý qua commit/PR. Quy trình yêu cầu đồng bộ SRS, sơ đồ/ADR, ma trận truy vết, code và test khi thay đổi hành vi; PR phải nêu lý do thay đổi.

\begin{table}[H]
\caption{Tình trạng tài liệu tại commit khảo sát}
\small
\begin{tabularx}{\linewidth}{L{0.23\linewidth}Y}
\toprule
\thead\textbf{Nhóm tài liệu} & \textbf{Kết quả và giới hạn} \\
\midrule
Task 1 & Có bối cảnh, phạm vi, SRS, use-case và sơ đồ trong \code{docs/requirements/}. \\
Task 2 & Có activity, sequence, class và state bổ sung; lưu cả nguồn \code{.puml} và bản \code{.svg}. Sự tồn tại của tệp chưa thay cho nghiệm thu nội dung. \\
Task 3 & \textbf{Còn thiếu} mô tả kiến trúc, ADR thực chất và deployment diagram; thư mục hiện có README hướng dẫn. \\
Ma trận truy vết & FR-01 và FR-02 có liên kết code/test; FR-03 mới có API chi tiết. Cột thiết kế vẫn cần hoàn thiện. \\
\bottomrule
\end{tabularx}
\end{table}

Nguồn: \gitfile{docs/traceability-matrix.md}{ma trận truy vết} và \gitfile{docs/architecture/README.md}{thư mục kiến trúc}. \textbf{Task 4.2 chưa đủ điều kiện nghiệm thu} vì chưa có đầy đủ tạo phẩm Task 1--3 theo thứ tự yêu cầu của đề.

\newpage
\section{Hiện thực MVP Sprint 1}

Backend sử dụng Catalog Controller/Service và Prisma để đọc PostgreSQL. Hai endpoint công khai phục vụ người dùng chưa đăng nhập:

\begin{table}[H]
\caption{API sản phẩm và hành vi phản hồi}
\small
\begin{tabularx}{\linewidth}{L{0.31\linewidth}Y}
\toprule
\thead\textbf{Endpoint} & \textbf{Kết quả} \\
\midrule
\code{GET /products} & \textbf{200}: mảng sản phẩm theo tên tăng dần; trả \code{[]} khi không có dữ liệu. \\
\code{GET /products/:id} & \textbf{200}: thông tin sản phẩm; \textbf{404} với mã \code{PRODUCT_NOT_FOUND} khi ID không tồn tại. \\
\bottomrule
\end{tabularx}
\end{table}

Sản phẩm có ID, tên, mô tả, giá nguyên VND, ảnh và tồn kho. Lỗi nội bộ trả \textbf{500} với thông báo chung, không đưa chi tiết lỗi database vào response.

\textbf{Giao diện menu.} Next.js lấy dữ liệu từ API; có loading, empty, error và nút thử lại. Thẻ món hiển thị ảnh, tên, giá, tồn kho; nhãn sắp hết áp dụng khi còn \textbf{1--5} sản phẩm, còn tồn \textbf{0} thì hiển thị hết hàng và vô hiệu hóa nút thêm vào giỏ. CSS bố trí \textbf{1/2/3 cột}, chuyển tại \textbf{560/900 px}; nút có chiều cao tối thiểu \textbf{44 px}. Đây là đặc điểm của mã nguồn, chưa phải kết quả đo trên thiết bị thật.

\textbf{Dữ liệu mẫu.} Migration khởi tạo bảng \code{products}; seed dùng \code{upsert} theo ID cho \textbf{8 món}. Ví dụ: Cappuccino \textbf{45.000đ}, tồn \textbf{30}; Bạc xỉu tồn \textbf{3}; Cold brew tồn \textbf{0}. Seed chạy lại có thể cập nhật giá/tồn kho về mẫu. Chưa có kiểm chứng tự động migrate/seed lặp trên database trống.

\begin{table}[H]
\caption{Phạm vi kiểm thử và kết quả local ngày 23/10/2026}
\small
\begin{tabularx}{\linewidth}{L{0.25\linewidth}Y C{0.10\linewidth}}
\toprule
\thead\textbf{Nhóm kiểm thử} & \textbf{Hành vi được kiểm tra} & \textbf{Đạt} \\
\midrule
API-T01--T02 & Danh sách và chi tiết sản phẩm trả 200, dữ liệu đúng. & 2/2 \\
API-T03 & ID không tồn tại trả 404 \code{PRODUCT_NOT_FOUND}. & 1/1 \\
API-T04 & Danh sách rỗng trả 200 và \code{[]}. & 1/1 \\
API-T05 & Lỗi DB trả 500, không lộ thông tin nội bộ. & 1/1 \\
\midrule
Web -- format & Định dạng giá VND. & 3/3 \\
Web -- stock & Nhãn còn hàng, sắp hết và hết hàng. & 3/3 \\
Web -- ProductCard & Thông tin món, trạng thái nút theo tồn kho. & 3/3 \\
Web -- MenuScreen & Loading, 8 thẻ món, empty, mất mạng/HTTP 503 và retry. & 4/4 \\
\midrule
\textbf{Tổng} & \textbf{API 5/5 và web 13/13; cả hai lệnh exit 0.} & \textbf{18/18} \\
\bottomrule
\end{tabularx}
\end{table}

{\small\raggedright
\textbf{Lệnh kiểm:} \code{npm.cmd --prefix apps/api test -- --runInBand}\\
\code{npm.cmd --prefix apps/web test}\par}
API test chạy HTTP qua NestJS với Prisma giả; web test mô phỏng fetch. \textbf{Chưa thay thế kiểm thử xuyên suốt với PostgreSQL thật trên trình duyệt.}

Nguồn: \gitfile{apps/api/test/products.e2e-spec.ts}{API test}, \gitfile{apps/web/src/components/menu/MenuScreen.test.tsx}{MenuScreen test} và \gitfile{apps/api/prisma/seed.ts}{seed}. Nút thêm vào giỏ hiện mới có hiển thị/disabled; chi tiết UI, giỏ hàng, xác thực, đơn hàng, thanh toán và PWA thuộc Sprint 2.

\newpage
\section{Continuous Integration và minh chứng}

\gitfile{.github/workflows/ci.yml}{Workflow CI} tự chạy khi mở/cập nhật PR vào \code{develop-Assignment2} và hỗ trợ chạy tay. Một job \code{ci} trên \code{ubuntu-latest} thực hiện \textbf{Install $\rightarrow$ Lint $\rightarrow$ Unit Test $\rightarrow$ Build}; dùng Node từ \code{.nvmrc}, quyền \code{contents: read}, không dùng \code{continue-on-error}. Test không cần database hoặc secret kết nối.

\begin{table}[H]
\caption{Đối chiếu hai lần CI đã lưu ngày 21/10/2026}
\small
\begin{tabularx}{\linewidth}{L{0.14\linewidth}Y L{0.18\linewidth}L{0.19\linewidth}}
\toprule
\thead\textbf{Bước} & \textbf{Nội dung} & \textbf{PR \#9} & \textbf{PR \#10 demo} \\
\midrule
Install & Cài hai app, generate Prisma Client. & Thành công & Thành công \\
Lint & ESLint API và web. & Thành công & Thành công \\
Unit Test & Jest API và Vitest web. & API 5/5; web 13/13 & API 5/5; web 12 đạt, 1 lỗi \\
Build & Web typecheck, NestJS build, Next.js build. & Thành công & \textbf{Bỏ qua} \\
\bottomrule
\end{tabularx}
\end{table}

\begin{itemize}
  \item \textbf{CI thành công:} \href{https://github.com/DuVinhLoi/BrewLite/pull/9}{PR \#9}, run \href{https://github.com/DuVinhLoi/BrewLite/actions/runs/39555707525}{39555707525}, commit \code{8a39ea9}; kết quả \code{success}. Có \gitfile{docs/evidence/ci/github-green.txt}{log xanh} và metadata trong repository.
  \item \textbf{CI phát hiện lỗi:} \href{https://github.com/DuVinhLoi/BrewLite/pull/10}{PR \#10}, run \href{https://github.com/DuVinhLoi/BrewLite/actions/runs/39555838160}{39555838160}; cố ý đặt kỳ vọng giá \textbf{0đ} thay vì \textbf{45.000đ}. Unit Test exit \textbf{1}, Build không chạy. Có \gitfile{docs/evidence/ci/github-red.txt}{log đỏ}; hồ sơ xác nhận PR đã đóng, không merge, test sai không vào nhánh Task 4.
\end{itemize}

\textbf{Điều kiện chặn merge còn thiếu.} Đề yêu cầu không merge khi lint/test/build lỗi. CI đỏ và Build bị bỏ qua mới chứng minh pipeline dừng; chưa chứng minh merge bị chặn. \gitfile{docs/evidence/ci/github-branch-protection.txt}{Hồ sơ GitHub ngày 21/10} ghi HTTP \textbf{403} do giới hạn gói repository private; tài khoản kiểm tra không có admin. Chưa có bằng chứng required check \code{ci} đã được bật và cưỡng chế thành công.

\section{Đánh giá kết quả và phần cần hoàn thiện}

\begin{table}[H]
\caption{Tổng hợp mức đáp ứng Task 4}
\small
\begin{tabularx}{\linewidth}{L{0.14\linewidth}Y Y}
\toprule
\thead\textbf{Hạng mục} & \textbf{Kết quả có minh chứng} & \textbf{Cần hoàn thiện hoặc xác minh} \\
\midrule
Task 4.1 & Cấu trúc repo, cấu hình mẫu, quy trình PR/commit/DoD, backlog 2 Sprint. & Phân công cá nhân, lịch Sprint; xác minh bảo vệ nhánh. \\
Task 4.2 & Tài liệu Task 1--2, lịch sử Git/PR, code/test trong ma trận. & Tạo phẩm Task 3 và liên kết thiết kế còn thiếu. \\
Task 4.3 & 2 API, menu, 8 món mẫu; local \textbf{18/18 test đạt}. & Kiểm thiết bị thật, luồng với DB thật, migrate/seed lặp trên DB trống. \\
Task 4.4 & CI tự chạy; có log xanh/đỏ thật, lỗi test khiến Build bỏ qua. & Required check \code{ci} và bằng chứng chặn merge khi lỗi. \\
\bottomrule
\end{tabularx}
\end{table}

\textbf{Kết luận:} Sprint 1 đã có MVP xem menu và CI kiểm tra tự động. \textbf{Chưa đủ điều kiện nghiệm thu toàn bộ Task 4} tại mốc khảo sát do thiếu tạo phẩm Task 3 và minh chứng chặn merge. Docker Compose, health check, CD, production, rollback và luồng đặt/thanh toán tiếp tục ở Task 5.

{\footnotesize\textbf{Phạm vi số liệu:} Đối chiếu \gitfile{docs/Assignment2-BrewLite.pdf}{đề A2 mục 6.1--6.4}; mã nguồn/tài liệu tại \code{0030c2a}; kết quả local 23/10 và hồ sơ CI 21/10/2026. Bản rút gọn không chạy lại test, lint, build hoặc migration.}
